\scanpage{123}
1) $\Phi$, pour la relation d'ordre induite par l'inclusion, est isomorphe à $\Phi_\tau$.

[2) \deleted{\unsure{Pour}} Pour tout $V\in\Phi$, on a $s_0\notin V$ et $V\unsure{\not\subset|\notin}H_\infty$ \uline{fibre par fibre}]

% A curved arrow connects condition 1 to the continuation below condition 2.
et pour $V\in\Phi$, la dimension combinatoire de $\Phi$ dans l'ens ordonné $\Phi$ (en prenant $-1$ comme dim du plus petit él) est \uline{égale} à la dim. projective de $V$. \deleted{\unsure{(donc dim $E=n$)}} Enfin $P\in\Phi$ (de dim $P=n$)

3) \deleted{\unsure{Tous}} Les éléments de $\Phi_0$ engendrent $P$ projectivement (fibre par fibre), et ceux de $\Phi_{n-1}$ engendrent $\check P$.
\deleted{Ceci implique que tt autom de $P$ qui invarie $\Phi$ est connu quand on connaît sa restriction à $\Phi$ comme \deleted{\unsure{autom}}\inserted{permutation} de l'ens ordonné $\Phi$.}

% margin note written sideways; arrow points to the position before condition 4
\inserted{Plus précisément, on veut que pour tout $V\subset W$, $V,W\in\Phi$, $W\ne V$ soit \deleted{\unsure{$\longrightarrow$}} le sous-espace projectif de $P$ engendré par les $V'\in\Phi$ tels que $V\subset V'\subset W$, $\dim V'=\dim V+1$ --- et \unsure{énoncé dual}.}

4) Tout automorphisme de l'ens ordonné $\Phi$ est induit par un autom. de $P$ invariant $\Phi$.
% archive sheet number: 122
